\documentclass[11pt,a4paper]{article}
\usepackage[margin=25mm]{geometry}
\usepackage{kotex}
\usepackage{hyperref}
\usepackage{xcolor}
\usepackage{tcolorbox}
\usepackage{booktabs}
\usepackage{graphicx}
\usepackage{amsmath}
\usepackage{amssymb}
\usepackage{listings}
\usepackage{setspace}
\usepackage{fancyhdr}
\usepackage{adjustbox}

% PDF 메타데이터 설정
\hypersetup{
  pdftitle={CSCI89 Lecture 11: 확률 모델과 생성 모델},
  pdfauthor={UIUC Student},
  pdfsubject={CSCI E-89: Introduction to Artificial Intelligence}
}

% 레이아웃 설정
\onehalfspacing % 줄 간격 1.5배
\setlength{\parskip}{0.6em} % 단락 사이 간격
\setlength{\parindent}{0pt} % 단락 들여쓰기 없음
\renewcommand{\arraystretch}{1.2} % 표 행 높이 조절

% 색상 정의
\definecolor{myblue}{HTML}{E0F2F7}
\definecolor{mygreen}{HTML}{E8F5E9}
\definecolor{myred}{HTML}{FFEBEE}
\definecolor{myyellow}{HTML}{FFFDE7}

% tcolorbox 스타일 정의
\tcbset{
  summarybox/.style={
    colback=myblue,
    colframe=blue!75!black,
    fonttitle=\bfseries,
    title style={fill=blue!20!white},
    arc=4mm, boxrule=0.8pt,
    left=2mm, right=2mm, top=2mm, bottom=2mm,
    boxsep=1mm,
    breakable
  },
  warningbox/.style={
    colback=myred,
    colframe=red!75!black,
    fonttitle=\bfseries,
    title style={fill=red!20!white},
    arc=4mm, boxrule=0.8pt,
    left=2mm, right=2mm, top=2mm, bottom=2mm,
    boxsep=1mm,
    breakable
  },
  examplebox/.style={
    colback=mygreen,
    colframe=green!75!black,
    fonttitle=\bfseries,
    title style={fill=green!20!white},
    arc=4mm, boxrule=0.8pt,
    left=2mm, right=2mm, top=2mm, bottom=2mm,
    boxsep=1mm,
    breakable
  },
  conceptbox/.style={
    colback=myyellow,
    colframe=orange!75!black,
    fonttitle=\bfseries,
    title style={fill=orange!20!white},
    arc=4mm, boxrule=0.8pt,
    left=2mm, right=2mm, top=2mm, bottom=2mm,
    boxsep=1mm,
    breakable
  }
}

% listings 설정 (코드 블록)
\lstset{
  basicstyle=\ttfamily\small, % 폰트 스타일 및 크기
  breaklines=true, % 긴 줄 자동 줄바꿈
  numbers=left, % 줄 번호 왼쪽에 표시
  numbersep=6pt, % 줄 번호와 코드 사이 간격
  frame=single, % 코드 블록 주위에 단일 프레임
  captionpos=b, % 캡션 위치 (아래)
  keywordstyle=\color{blue}, % 키워드 색상
  stringstyle=\color{red}, % 문자열 색상
  commentstyle=\color{green!50!black}, % 주석 색상
  showstringspaces=false, % 문자열 내 공백 표시 안 함
  tabsize=2, % 탭 크기
  columns=fullflexible, % 유연한 열 너비
  keepspaces=true, % 공백 유지
  xleftmargin=1em, xrightmargin=1em % 좌우 여백
}

% 헤더/푸터 설정
\pagestyle{fancy}
\fancyhf{} % 모든 헤더/푸터 초기화
\fancyhead[L]{CSCI89 Lecture 11: 확률 모델과 생성 모델} % 왼쪽 헤더
\fancyhead[R]{\today} % 오른쪽 헤더
\fancyfoot[C]{\thepage} % 중앙 푸터에 페이지 번호

\begin{document}

\begin{tcolorbox}[summarybox, title=개요]
본 문서는 CSCI89 Lecture 11 강의 내용을 바탕으로 인공지능 분야의 핵심 확률 모델과 생성 모델을 다룹니다.
강의는 퀴즈 10 복습으로 시작하여, 마르코프 체인(Markov Chains)의 기본 개념을 소개하고, 이를 확장한 은닉 마르코프 모델(Hidden Markov Models, HMM)을 설명합니다.
이어서 HMM의 한계를 극복하기 위한 조건부 무작위 필드(Conditional Random Fields, CRF)와, CRF의 수동 특징 공학 문제를 해결하는 양방향 LSTM-CRF(Bi-directional LSTM-CRF) 결합 모델을 학습합니다.
마지막으로, 잠재 공간의 구조화 문제를 해결하는 변분 오토인코더(Variational Autoencoders, VAE)와, 비지도 학습을 통해 현실적인 데이터를 생성하는 생성적 적대 신경망(Generative Adversarial Networks, GAN)의 원리를 탐구합니다.
각 모델의 필요성, 작동 원리, 장단점, 그리고 실제 적용 사례를 비전공자도 이해하기 쉽게 설명하는 데 중점을 둡니다.
\end{tcolorbox}

\newpage
\section*{용어 정리}

다음 표는 본 강의에서 다루는 주요 용어들을 쉽게 설명하고 원어와 함께 제시합니다.

\begin{table}[h!]
\centering
\caption{주요 용어 정리}
\label{tab:glossary}
\begin{adjustbox}{width=\textwidth}
\begin{tabular}{lllc}
\toprule
\textbf{용어} & \textbf{쉬운 설명} & \textbf{원어} & \textbf{비고} \\
\midrule
개체명 인식 & 텍스트에서 사람 이름, 장소, 조직 등 특정 개체를 찾아 분류하는 기술 & Named Entity Recognition (NER) & NLP의 핵심 과제 \\
마르코프 체인 & 미래 상태가 현재 상태에만 의존하는 확률적 과정 & Markov Chain & HMM의 기반 \\
마르코프 속성 & 과거 정보가 주어져도 미래는 현재 상태에만 의존한다는 가정 & Markov Property & 많은 확률 모델의 기본 가정 \\
은닉 마르코프 모델 & 관찰 불가능한(은닉) 상태와 관찰 가능한 상태가 연결된 확률 모델 & Hidden Markov Model (HMM) & 시퀀스 데이터 모델링에 사용 \\
전이 확률 & 한 은닉 상태에서 다른 은닉 상태로 이동할 확률 & Transition Probability & HMM의 핵심 파라미터 \\
방출 확률 & 특정 은닉 상태에서 특정 관찰 가능한 상태가 나타날 확률 & Emission Probability & HMM의 핵심 파라미터 \\
기대-최대화 알고리즘 & 은닉 변수가 있는 모델의 파라미터를 추정하는 반복적 방법 & Expectation-Maximization (EM) Algorithm & HMM 훈련에 사용 \\
조건부 무작위 필드 & 입력 시퀀스 전체에 대한 조건부 확률을 모델링하는 확률 모델 & Conditional Random Field (CRF) & HMM의 독립성 가정 완화 \\
특징 함수 & 입력 데이터와 상태 간의 특정 패턴을 포착하는 수동 정의 함수 & Feature Function & CRF의 핵심 구성 요소 \\
양방향 LSTM & 과거와 미래 정보를 모두 사용하여 시퀀스 데이터를 처리하는 신경망 & Bi-directional Long Short-Term Memory (Bi-LSTM) & 문맥 이해에 효과적 \\
변분 오토인코더 & 잠재 공간을 구조화하여 의미 있는 데이터 생성을 가능하게 하는 모델 & Variational Autoencoder (VAE) & 생성 모델의 일종 \\
쿨백-라이블러 발산 & 두 확률 분포 간의 차이를 측정하는 방법 & Kullback-Leibler (KL) Divergence & VAE 손실 함수에 사용 \\
생성적 적대 신경망 & 생성자와 판별자가 서로 경쟁하며 학습하는 생성 모델 & Generative Adversarial Network (GAN) & 현실적인 데이터 생성에 탁월 \\
생성자 & 실제와 유사한 데이터를 생성하는 신경망 & Generator & GAN의 한 부분 \\
판별자 & 입력 데이터가 실제인지 생성된 것인지 구별하는 신경망 & Discriminator & GAN의 한 부분 \\
\bottomrule
\end{tabular}
\end{adjustbox}
\end{table}

\newpage
\section*{핵심 개념·원리}

\subsection{퀴즈 10 복습}

강의는 지난 퀴즈 10의 주요 질문들을 복습하며 시작했습니다. 이는 자연어 처리(NLP) 분야의 핵심 개념들을 다시 한번 상기시키는 기회가 됩니다.

\begin{tcolorbox}[conceptbox, title=퀴즈 10 주요 내용]
\begin{enumerate}
    \item \textbf{개체명 인식(Named Entity Recognition, NER)의 주요 목적:} 텍스트 내에서 이름이 지정된 개체(예: 사람, 장소, 조직)를 탐지하고 분류하는 것입니다. (정답: B)
    \item \textbf{규칙 기반 개체명 인식 접근 방식:} 미리 정의된 언어 규칙과 패턴을 사용하여 텍스트에서 개체를 체계적으로 식별합니다. (정답: C)
    \item \textbf{spaCy 라이브러리의 역할:} 다양한 언어에 대해 미리 훈련된 개체명 인식 모델을 제공합니다. (정답: C)
    \begin{itemize}
        \item \textbf{학생 질문:} spaCy가 파이프라인을 시각화하는 기능(B)은 왜 정답이 아닌가요?
        \item \textbf{교수 답변:} spaCy는 시각화 도구를 제공하지만, 이는 정적인 이미지나 시각화일 뿐 상호작용 가능한 그래픽 사용자 인터페이스(GUI)는 아닙니다. 즉, 클릭하여 확장하거나 변경할 수 없습니다. 따라서 '시각화'는 하지만 'GUI'는 아닙니다.
    \end{itemize}
    \item \textbf{규칙 기반 개체명 인식 시스템의 확장성 문제:} 규칙 세트가 증가할수록 시스템이 점점 더 복잡해지고 효과적으로 관리하기 어려워집니다. (정답: C)
    \item \textbf{통계적 개체명 인식 방법의 주요 단점:} 훈련을 위해 대규모의 잘 주석된 데이터 세트가 필요하며, 이는 확보하고 유지하는 데 많은 자원이 소모됩니다. (정답: A)
\end{enumerate}
\end{tcolorbox}

\subsection{마르코프 체인 (Markov Chains)}

\subsubsection{개념 및 정의}
마르코프 체인(Markov Chain)은 미래의 상태가 오직 현재의 상태에만 의존하고, 과거의 상태에는 독립적인 확률적 과정을 의미합니다. 이는 "마르코프 속성(Markov Property)"이라는 핵심 가정을 기반으로 합니다.

\begin{tcolorbox}[conceptbox, title=마르코프 속성 (Markov Property)]
\textbf{한 줄 요약:} 미래는 현재에만 달려있고, 과거는 중요하지 않다.

\textbf{직관적 예시:} 날씨를 예측할 때, 내일 날씨는 오늘 날씨에만 영향을 받고, 어제나 그제 날씨는 직접적인 영향을 주지 않는다고 가정하는 것과 같습니다.

\textbf{기술적 설명:} 시점 $t+1$의 상태 $X_{t+1}$의 확률은 시점 $t$의 상태 $X_t$가 주어졌을 때, 그 이전의 모든 상태들($X_{t-1}, X_{t-2}, \dots$)이 주어지는 것과 동일합니다. 즉, 수학적으로 다음과 같이 표현됩니다.
\[ P(X_{t+1} = x_{t+1} | X_t = x_t, X_{t-1} = x_{t-1}, \dots, X_1 = x_1) = P(X_{t+1} = x_{t+1} | X_t = x_t) \]
여기서 $x_t$는 시점 $t$에서의 특정 상태를 나타냅니다.
\end{tcolorbox}

\subsubsection{작동 원리 및 예시}
마르코프 체인은 일련의 상태(state)와 상태 간의 전이 확률(transition probability)로 구성됩니다.

\begin{examplebox}[title=마르코프 체인 예시: 간단한 어휘 모델]
세 가지 단어로 구성된 어휘(상태 1, 상태 2, 상태 3)가 있다고 가정해 봅시다.
현재 상태가 '상태 1'일 때, 다음 단어가 '상태 1', '상태 2', '상태 3' 중 어느 것이 될지 확률로 결정됩니다.

\begin{itemize}
    \item $P_{11}$: 상태 1에서 상태 1로 전이할 확률
    \item $P_{12}$: 상태 1에서 상태 2로 전이할 확률
    \item $P_{13}$: 상태 1에서 상태 3으로 전이할 확률
\end{itemize}
이때, 한 상태에서 다음 상태로 전이할 모든 확률의 합은 1이 되어야 합니다.
\[ P_{11} + P_{12} + P_{13} = 1 \]
모든 전이 확률 $P_{ij}$는 0보다 크거나 같아야 합니다.
이러한 방식으로, 현재 단어(상태)가 다음 단어(상태)에 영향을 미치는 언어 모델을 만들 수 있습니다.
\end{examplebox}

\subsubsection{장점과 한계}
\begin{itemize}
    \item \textbf{장점:} 개념이 간단하고 구현하기 쉽습니다. 특정 시스템의 동역학을 모델링하는 데 유용합니다.
    \item \textbf{한계:} 마르코프 속성이라는 강력한 가정이 현실 세계의 복잡한 의존성을 포착하기 어렵게 만듭니다. 예를 들어, 문장에서 다음 단어는 바로 직전 단어뿐만 아니라 문맥 전체에 영향을 받을 수 있습니다.
\end{itemize}

\begin{tcolorbox}[warningbox, title=마르코프 속성의 한계 극복]
마르코프 속성이 너무 강력하다고 느껴질 때, "현재 상태"의 정의를 확장하여 이 한계를 완화할 수 있습니다.
\textbf{예시:}
\begin{itemize}
    \item \textbf{자율 주행:} 현재 차량의 위치만으로 미래를 예측하는 대신, 현재 위치, 속도, 방향, 주변 객체의 움직임 등 여러 스냅샷(짧은 비디오)을 묶어 "현재 상태"로 정의할 수 있습니다.
    \item \textbf{자연어 처리:} 다음 단어가 이전 단어 하나에만 의존하는 것이 아니라, 이전 두 단어(2-gram)나 세 단어(3-gram)에 의존하도록 "상태"를 정의할 수 있습니다. 이렇게 하면 더 많은 과거 정보를 "현재 상태"에 포함시켜 마르코프 속성을 유지하면서도 더 현실적인 모델을 만들 수 있습니다.
\end{itemize}
\end{tcolorbox}

\subsection{은닉 마르코프 모델 (Hidden Markov Models, HMM)}

\subsubsection{개념 및 필요성}
마르코프 체인은 모든 상태가 관찰 가능하다는 가정을 합니다. 하지만 현실에서는 우리가 직접 관찰할 수 없는 "은닉된" 요인들이 시스템의 동역학에 영향을 미치는 경우가 많습니다. 은닉 마르코프 모델(HMM)은 이러한 은닉 상태(hidden states)의 개념을 도입하여 마르코프 체인을 확장한 모델입니다.

\begin{tcolorbox}[conceptbox, title=은닉 마르코프 모델 (HMM)의 핵심 아이디어]
\textbf{한 줄 요약:} 우리가 보는 것은 빙산의 일각일 뿐, 그 아래에는 보이지 않는 의도가 숨어있다.

\textbf{직관적 예시:} 우리가 말을 할 때, 실제로 발화되는 단어(관찰 가능한 상태) 뒤에는 우리가 말하고자 하는 "의도"나 "품사"와 같은 은닉된 상태가 있다고 생각할 수 있습니다. 예를 들어, "고양이"라는 단어를 말할 때, 그 단어는 '명사'라는 은닉 상태에서 생성된 것입니다.

\textbf{기술적 설명:} HMM은 두 가지 종류의 상태 시퀀스를 가집니다.
\begin{itemize}
    \item \textbf{은닉 상태 (Hidden States, $Y$):} 직접 관찰할 수 없지만, 시스템의 내부 동역학을 나타냅니다. 이 은닉 상태들은 마르코프 속성을 따릅니다 (즉, $Y_t$는 $Y_{t-1}$에만 의존).
    \item \textbf{관찰 가능한 상태 (Observable States, $X$):} 각 은닉 상태에서 생성되는, 우리가 실제로 관찰할 수 있는 데이터입니다.
\end{itemize}
\end{tcolorbox}

\subsubsection{HMM의 구조}
HMM은 은닉 상태 시퀀스($Y_1, Y_2, \dots, Y_T$)와 관찰 가능한 상태 시퀀스($X_1, X_2, \dots, X_T$)로 구성됩니다. 각 $Y_t$는 $Y_{t-1}$에 따라 전이하고, 각 $X_t$는 $Y_t$에 따라 생성됩니다.

\begin{center}
\includegraphics[width=0.8\textwidth]{hmm_structure.png} % 이미지 파일이 없으므로 주석 처리합니다. 실제 사용 시 이미지 경로를 지정하세요.
\end{center}
\textit{HMM 구조의 개념도: 은닉 상태(Y)는 서로 전이하며, 각 은닉 상태에서 관찰 가능한 상태(X)가 생성됩니다.}

\subsubsection{HMM의 파라미터}
HMM을 정의하는 세 가지 주요 파라미터 집합이 있습니다.
\begin{enumerate}
    \item \textbf{전이 확률 (Transition Probabilities):} 한 은닉 상태에서 다른 은닉 상태로 이동할 확률입니다.
    \[ A_{ij} = P(Y_t = j | Y_{t-1} = i) \]
    \item \textbf{방출 확률 (Emission Probabilities):} 특정 은닉 상태에서 특정 관찰 가능한 상태가 생성될 확률입니다.
    \[ B_j(k) = P(X_t = k | Y_t = j) \]
    \item \textbf{초기 상태 확률 (Initial State Probabilities):} 시퀀스의 시작 시점에 각 은닉 상태에 있을 확률입니다.
    \[ \pi_i = P(Y_1 = i) \]
\end{enumerate}

\subsubsection{HMM의 훈련 및 적용}
HMM은 관찰된 데이터($X$)를 기반으로 이 세 가지 파라미터를 추정하여 훈련됩니다.

\begin{itemize}
    \item \textbf{훈련 (Training):}
    \begin{itemize}
        \item \textbf{은닉 상태($Y$)가 관찰 가능한 경우:} 만약 훈련 데이터에 은닉 상태(예: 품사 태그)가 모두 주석되어 있다면, 단순히 빈도수를 세어 최대 우도(Maximum Likelihood) 추정을 통해 파라미터를 직접 계산할 수 있습니다.
        \item \textbf{은닉 상태($Y$)가 관찰 불가능한 경우:} 은닉 상태를 알 수 없을 때는 기대-최대화(Expectation-Maximization, EM) 알고리즘을 사용하여 파라미터를 추정합니다. EM 알고리즘은 현재 파라미터로 은닉 상태를 추정(E-단계)하고, 추정된 은닉 상태를 바탕으로 파라미터를 업데이트(M-단계)하는 과정을 반복합니다.
    \end{itemize}
    \item \textbf{적용 (Applications):}
    \begin{itemize}
        \item \textbf{품사 태깅 (Part-of-Speech Tagging):} $Y$는 품사(명사, 동사 등), $X$는 실제 단어가 됩니다. 훈련 후, 새로운 문장($X$)이 주어졌을 때 가장 가능성 높은 품사 시퀀스($Y$)를 예측합니다.
        \item \textbf{음성 인식 (Speech Recognition):} $Y$는 실제 텍스트(단어), $X$는 음성 신호(소리 특징)가 됩니다. 음성 데이터를 텍스트로 변환하는 데 사용됩니다.
    \end{itemize}
\end{itemize}

\subsubsection{HMM의 한계}
HMM은 마르코프 체인과 마찬가지로 강력한 독립성 가정을 가집니다.
\begin{itemize}
    \item \textbf{단순화된 의존성:} 현재 은닉 상태($Y_t$)는 오직 직전 은닉 상태($Y_{t-1}$)에만 의존합니다. 또한, 관찰 가능한 상태($X_t$)는 오직 현재 은닉 상태($Y_t$)에만 의존합니다. 이는 문맥 전체나 장거리 의존성을 포착하기 어렵게 만듭니다.
    \item \textbf{예시:} 품사 태깅에서 "명사"라는 은닉 상태가 다음 품사에 영향을 미치지만, 그 명사가 어떤 "단어"인지(예: "고양이"인지 "책"인지)는 다음 품사 결정에 직접적인 영향을 주지 않는다고 가정합니다. 이는 비현실적일 수 있습니다.
\end{itemize}

\subsection{조건부 무작위 필드 (Conditional Random Fields, CRF)}

\subsubsection{개념 및 필요성}
HMM의 가장 큰 한계는 강력한 독립성 가정입니다. 특히, 관찰 가능한 상태($X_t$)가 오직 현재 은닉 상태($Y_t$)에만 의존하고, 다른 관찰 가능한 상태($X_{t-1}, X_{t+1}$ 등)에는 독립적이라는 가정은 현실 세계의 복잡한 패턴을 모델링하기 어렵게 만듭니다. 조건부 무작위 필드(CRF)는 이러한 HMM의 한계를 극복하기 위해 도입되었습니다.

\begin{tcolorbox}[conceptbox, title=조건부 무작위 필드 (CRF)의 핵심 아이디어]
\textbf{한 줄 요약:} HMM보다 훨씬 유연하게, 문맥 전체를 고려하여 예측한다.

\textbf{직관적 예시:} HMM이 "이 단어가 명사라면 다음 단어는 동사일 확률이 높다"고 예측한다면, CRF는 "이 단어가 'Mr.'이고 다음 단어가 대문자로 시작하는 명사라면, 이 단어는 사람 이름일 확률이 높다"와 같이 훨씬 더 복잡하고 문맥적인 규칙을 고려하여 예측합니다.

\textbf{기술적 설명:} CRF는 입력 시퀀스 $X$가 주어졌을 때, 출력 시퀀스 $Y$의 조건부 확률 $P(Y|X)$를 직접 모델링합니다. HMM과 달리, CRF는 $Y$와 $X$ 사이의 복잡한 비선형적 관계를 포착할 수 있으며, $Y_t$가 $X$ 시퀀스 전체($X_1, \dots, X_T$)에 의존할 수 있도록 합니다. 이는 방향성이 없는 그래프 모델(undirected graphical model)로 표현됩니다.
\end{tcolorbox}

\subsubsection{CRF의 수학적 표현}
CRF의 조건부 확률은 다음과 같이 정의됩니다.
\[ P(Y|X) = \frac{1}{Z(X)} \exp\left(\sum_k \lambda_k f_k(Y_t, Y_{t-1}, X)\right) \]
여기서 각 요소는 다음과 같습니다.
\begin{itemize}
    \item $Y$: 은닉 상태(출력) 시퀀스
    \item $X$: 관찰 가능한(입력) 시퀀스 (벡터 $X$는 $X_1, \dots, X_T$ 전체를 의미)
    \item $Z(X)$: 정규화 상수(Normalization Factor). 모든 가능한 $Y$ 시퀀스에 대한 $\exp(\dots)$ 값의 합으로, 확률의 총합이 1이 되도록 만듭니다. (소프트맥스(softmax)의 분모와 유사)
    \item $\lambda_k$: 각 특징 함수 $f_k$에 대한 가중치(weight) 파라미터. 이 값들은 훈련 데이터를 통해 추정됩니다.
    \item $f_k(Y_t, Y_{t-1}, X)$: \textbf{특징 함수 (Feature Function)}. CRF의 핵심 요소로, 입력 시퀀스 $X$의 특정 부분, 현재 은닉 상태 $Y_t$, 그리고 이전 은닉 상태 $Y_{t-1}$ 간의 관계를 포착하는 함수입니다.
\end{itemize}

\subsubsection{특징 함수 ($f_k$)의 역할과 예시}
특징 함수는 CRF 모델에 문맥 정보를 주입하는 역할을 합니다. 이 함수들은 일반적으로 수동으로 설계되며, 특정 패턴이 나타날 때 1을 반환하고 그렇지 않으면 0을 반환하는 이진(binary) 지시 함수(indicator function) 형태를 가집니다.

\begin{examplebox}[title=CRF 특징 함수 예시: 사람 이름 인식]
사람 이름을 인식하는 NER 태스크를 가정해 봅시다. 다음과 같은 특징 함수를 만들 수 있습니다.

\begin{itemize}
    \item \textbf{특징 함수 1 ($f_1$):}
    "만약 $X_{t-1}$이 'Mr.' 또는 'Dr.'이고, $X_t$가 대문자로 시작하는 단어이며, $Y_{t-1}$이 명사이고 $Y_t$도 명사라면, $f_1 = 1$."
    (예: "Mr. John"에서 'John'이 사람 이름일 가능성을 높임)

    \item \textbf{특징 함수 2 ($f_2$):}
    "만약 $X_{t-2}$가 'Ms.'이고, $X_{t-1}$과 $X_t$ 모두 대문자로 시작하는 단어이며, $Y_{t-2}, Y_{t-1}, Y_t$가 모두 명사라면, $f_2 = 1$."
    (예: "Ms. Jane Doe"에서 'Jane Doe'가 사람 이름일 가능성을 높임)
\end{itemize}
이처럼 특징 함수는 $X$ 시퀀스 전체의 어떤 부분과 $Y$ 시퀀스의 특정 부분 간의 복잡한 관계를 정의할 수 있습니다. $\lambda_k$는 이러한 규칙들의 중요도를 학습합니다.
\end{examplebox}

\subsubsection{CRF의 장점과 한계}
\begin{itemize}
    \item \textbf{장점:}
    \begin{itemize}
        \item \textbf{유연성:} HMM의 강력한 독립성 가정을 완화하여, 입력 시퀀스 $X$의 모든 부분에 의존하는 특징 함수를 정의할 수 있습니다. 이는 장거리 의존성(long-range dependencies)과 풍부한 문맥 정보를 포착할 수 있게 합니다.
        \item \textbf{문맥 인식:} 과거뿐만 아니라 미래의 관찰 값($X$)까지 고려하여 현재 상태($Y_t$)를 예측할 수 있습니다.
        \item \textbf{강력한 성능:} 품사 태깅, 개체명 인식 등 시퀀스 레이블링(sequence labeling) 문제에서 HMM보다 뛰어난 성능을 보입니다.
    \end{itemize}
    \item \textbf{한계:}
    \begin{itemize}
        \item \textbf{집중적인 특징 공학 (Intensive Feature Engineering):} 특징 함수 $f_k$를 수동으로 설계해야 합니다. 이는 많은 시간, 도메인 지식, 그리고 노력을 요구하는 매우 어려운 작업입니다. 어떤 특징이 중요한지, 어떻게 조합해야 할지 결정하는 것이 모델의 성능에 결정적인 영향을 미칩니다.
        \item \textbf{높은 계산 비용:} $Z(X)$를 계산하는 과정이 복잡하고 계산 비용이 높을 수 있습니다.
    \end{itemize}
\end{itemize}

\subsection{양방향 LSTM과 CRF 결합 (Bi-directional LSTM-CRF)}

\subsubsection{개념 및 필요성}
CRF는 강력한 모델이지만, 수동으로 특징 함수를 설계해야 하는 "특징 공학(Feature Engineering)"이라는 큰 단점을 가지고 있습니다. 신경망, 특히 순환 신경망(Recurrent Neural Networks, RNN)의 일종인 LSTM(Long Short-Term Memory)은 시퀀스 데이터에서 자동으로 특징을 학습하는 데 탁월합니다. 양방향 LSTM(Bi-directional LSTM, Bi-LSTM)은 과거와 미래의 문맥 정보를 모두 포착할 수 있습니다.

Bi-LSTM과 CRF를 결합한 모델은 Bi-LSTM이 자동으로 풍부한 문맥 특징을 학습하고, CRF 레이어가 이 특징들을 바탕으로 시퀀스 레이블링의 제약 조건(예: 명사 뒤에 동사가 올 수 있지만, 명사 뒤에 명사가 오는 것은 드물다)을 모델링하여 최적의 레이블 시퀀스를 예측하도록 합니다.

\begin{tcolorbox}[conceptbox, title=Bi-LSTM-CRF의 핵심 아이디어]
\textbf{한 줄 요약:} 신경망이 똑똑하게 특징을 뽑아주고, CRF가 그 특징들을 활용해 시퀀스 규칙에 맞춰 최종 답을 내놓는다.

\textbf{직관적 예시:}
\begin{itemize}
    \item \textbf{Bi-LSTM:} 문장을 읽으면서 각 단어의 의미와 문맥을 파악합니다. "Mr. John Smith"라는 문장에서 'John'이라는 단어를 볼 때, Bi-LSTM은 'Mr.'(과거 문맥)와 'Smith'(미래 문맥)를 모두 고려하여 'John'이 사람 이름일 가능성이 높다는 특징을 추출합니다.
    \item \textbf{CRF:} Bi-LSTM이 추출한 특징들을 바탕으로, "사람 이름 뒤에는 또 다른 사람 이름이 올 수 있지만, 사람 이름 뒤에 동사가 오는 것은 어색하다"와 같은 시퀀스 규칙을 적용하여 최종적으로 가장 자연스러운 레이블 시퀀스("B-PER I-PER I-PER" 등)를 결정합니다.
\end{itemize}

\textbf{기술적 설명:} Bi-LSTM-CRF 모델은 다음과 같은 아키텍처를 가집니다.
\begin{enumerate}
    \item \textbf{임베딩 레이어 (Embedding Layer):} 입력 단어 시퀀스($X$)를 밀집 벡터(dense vector) 표현으로 변환합니다.
    \item \textbf{양방향 LSTM 레이어 (Bi-directional LSTM Layer):} 임베딩된 시퀀스를 입력받아 각 시점($t$)에서 과거와 미래의 문맥 정보를 모두 포함하는 은닉 상태 벡터($H_t$)를 출력합니다.
    \item \textbf{선형 변환 레이어 (Linear Transformation Layer):} Bi-LSTM의 은닉 상태 $H_t$를 입력받아 "방출 점수(Emission Scores)" 벡터 $E_t$를 생성합니다. 이 $E_t$는 각 시점 $t$에서 가능한 모든 출력 레이블(예: 품사)에 대한 점수를 나타냅니다.
    \[ E_t = W H_t + B \]
    여기서 $W$는 가중치 행렬, $B$는 편향 벡터입니다. $E_t$는 전통적인 CRF의 특징 함수 역할을 대체합니다.
    \item \textbf{CRF 레이어 (CRF Layer):} 선형 변환 레이어에서 나온 방출 점수 $E_t$를 입력으로 받아, 시퀀스 전체에 대한 제약 조건(예: 레이블 간의 전이 규칙)을 고려하여 가장 가능성 높은 출력 레이블 시퀀스($Y$)를 예측합니다.
\end{enumerate}
\end{tcolorbox}

\subsubsection{훈련 및 적용}
Bi-LSTM-CRF 모델은 종단 간(end-to-end)으로 훈련됩니다. 즉, 임베딩 레이어부터 CRF 레이어까지 모든 파라미터가 하나의 손실 함수(예: 음의 로그 우도, Negative Log-Likelihood)를 최소화하는 방향으로 동시에 학습됩니다.

\begin{itemize}
    \item \textbf{훈련:} 훈련 데이터($X$와 $Y$ 레이블 시퀀스)를 사용하여 모델의 모든 가중치(임베딩, LSTM, 선형 변환, CRF 전이 점수)를 최적화합니다.
    \item \textbf{적용:} 개체명 인식(NER), 품사 태깅(POS Tagging) 등 시퀀스 레이블링 문제에서 뛰어난 성능을 발휘합니다.
\item \textbf{장점:}
    \begin{itemize}
        \item \textbf{자동 특징 학습:} 수동 특징 공학의 필요성을 없애고, Bi-LSTM이 자동으로 문맥적 특징을 학습합니다.
        \item \textbf{문맥 포착:} 양방향 LSTM을 통해 과거와 미래의 문맥을 모두 효과적으로 활용합니다.
        \item \textbf{시퀀스 제약 조건 모델링:} CRF 레이어가 출력 레이블 시퀀스 간의 유효한 전이 규칙을 학습하여 비현실적인 레이블 조합을 방지합니다.
    \end{itemize}
    \item \textbf{단점:}
    \begin{itemize}
        \item \textbf{높은 계산 비용:} LSTM과 CRF의 결합으로 인해 훈련 및 추론 시간이 길고, 많은 컴퓨팅 자원을 요구합니다.
        \item \textbf{모델 튜닝의 복잡성:} 많은 하이퍼파라미터(LSTM의 은닉 차원, 임베딩 차원 등)를 튜닝해야 합니다.
    \end{itemize}
\end{itemize}

\subsubsection{Bi-LSTM-CRF 코드 구조 예시}
다음은 Bi-LSTM-CRF 모델의 파이토치(PyTorch)와 유사한 구현 구조를 보여줍니다.

\begin{lstlisting}[language=Python, caption={Bi-LSTM-CRF 모델의 간략한 PyTorch 유사 코드 구조}, label={lst:bilstm_crf_code}]
import torch
import torch.nn as nn
from torchcrf import CRF # CRF 레이어는 외부 라이브러리 사용

class BiLSTM_CRF(nn.Module):
    def __init__(self, vocab_size, tag_to_ix, embedding_dim, hidden_dim):
        super(BiLSTM_CRF, self).__init__()
        self.embedding_dim = embedding_dim
        self.hidden_dim = hidden_dim
        self.vocab_size = vocab_size
        self.tag_to_ix = tag_to_ix
        self.tagset_size = len(tag_to_ix)

        # 1. 임베딩 레이어: 단어를 벡터로 변환
        self.word_embeds = nn.Embedding(vocab_size, embedding_dim)

        # 2. 양방향 LSTM 레이어: 문맥 특징 학습
        # bidirectional=True로 양방향 LSTM 설정
        # hidden_dim은 단방향 LSTM의 은닉 차원. 양방향이므로 실제 출력은 2 * hidden_dim
        self.lstm = nn.LSTM(embedding_dim, hidden_dim // 2,
                            num_layers=1, bidirectional=True)

        # 3. 선형 변환 레이어: LSTM 출력을 태그 점수(emission scores)로 변환
        # LSTM 출력 (2 * hidden_dim)을 태그셋 크기(tagset_size)로 매핑
        self.hidden2tag = nn.Linear(hidden_dim, self.tagset_size)

        # 4. CRF 레이어: 태그 점수를 바탕으로 시퀀스 제약 조건 모델링
        self.crf = CRF(self.tagset_size, batch_first=True)

    def _get_lstm_features(self, sentence):
        # 임베딩
        embeds = self.word_embeds(sentence)
        # LSTM 입력 형태 (seq_len, batch_size, input_size)
        embeds = embeds.permute(1, 0, 2) # (batch_size, seq_len, embedding_dim) -> (seq_len, batch_size, embedding_dim)
        
        # LSTM 통과
        lstm_out, _ = self.lstm(embeds)
        # LSTM 출력 형태 (seq_len, batch_size, hidden_dim * num_directions)
        
        # 선형 변환으로 태그 점수(emission scores) 계산
        lstm_feats = self.hidden2tag(lstm_out)
        # 출력 형태 (seq_len, batch_size, tagset_size)
        
        return lstm_feats.permute(1, 0, 2) # (batch_size, seq_len, tagset_size)

    def forward(self, sentence, tags):
        # LSTM 특징(emission scores) 계산
        lstm_feats = self._get_lstm_features(sentence)
        
        # CRF 손실 계산 (음의 로그 우도)
        # tags는 실제 레이블 시퀀스
        loss = -self.crf(lstm_feats, tags, reduction='mean')
        return loss

    def decode(self, sentence):
        # LSTM 특징(emission scores) 계산
        lstm_feats = self._get_lstm_features(sentence)
        
        # CRF 디코딩 (가장 가능성 높은 레이블 시퀀스 예측)
        # viterbi_decode는 CRF에서 최적의 경로를 찾는 알고리즘
        best_tags = self.crf.decode(lstm_feats)
        return best_tags

# 예시 사용법 (훈련 데이터 준비)
# vocab_size: 어휘 사전 크기
# tag_to_ix: 태그(품사)를 인덱스로 매핑하는 딕셔너리 (예: {'NOUN': 0, 'VERB': 1, ...})
# embedding_dim: 임베딩 벡터 차원 (예: 100)
# hidden_dim: LSTM 은닉 상태 차원 (예: 256)

# model = BiLSTM_CRF(vocab_size, tag_to_ix, embedding_dim, hidden_dim)
# optimizer = optim.SGD(model.parameters(), lr=0.01)

# for epoch in range(num_epochs):
#     for sentence, tags in training_data:
#         # 훈련 데이터는 이미 인덱스로 변환되어 있어야 함
#         # sentence: (batch_size, seq_len)
#         # tags: (batch_size, seq_len)
#         
#         optimizer.zero_grad()
#         loss = model(sentence, tags)
#         loss.backward()
#         optimizer.step()

# # 예측 예시
# # test_sentence: (1, seq_len)
# # predicted_tags = model.decode(test_sentence)
\end{lstlisting}

\begin{tcolorbox}[summarybox, title=코드 설명 요약]
\begin{itemize}
    \item \textbf{입력:} 문장(단어 인덱스 시퀀스)과 해당 문장의 실제 태그(레이블 인덱스 시퀀스)를 받습니다.
    \item \textbf{임베딩 레이어 (\texttt{word\_embeds}):} 각 단어 인덱스를 고차원 벡터로 변환합니다.
    \item \textbf{LSTM 레이어 (\texttt{lstm}):} 임베딩된 단어 시퀀스를 처리하여 각 단어에 대한 문맥 정보를 담은 은닉 상태 벡터를 생성합니다. \texttt{bidirectional=True}로 설정되어 과거와 미래 문맥을 모두 고려합니다.
    \item \textbf{선형 레이어 (\texttt{hidden2tag}):} LSTM의 은닉 상태를 각 가능한 태그(레이블)에 대한 "점수(score)" 또는 "방출 점수(emission score)"로 변환합니다. 이 점수들은 CRF 레이어의 입력이 됩니다.
    \item \textbf{CRF 레이어 (\texttt{crf}):} 방출 점수와 실제 태그 시퀀스를 사용하여 손실을 계산합니다. 이 손실은 태그 시퀀스 간의 전이 규칙을 학습하는 데 사용됩니다. \texttt{decode} 메서드는 훈련된 CRF를 사용하여 새로운 문장에 대한 최적의 태그 시퀀스를 예측합니다.
\end{itemize}
\end{tcolorbox}

\newpage
\subsection{변분 오토인코더 (Variational Autoencoders, VAE)}

\subsubsection{오토인코더 (Autoencoders, AE)의 한계}
변분 오토인코더(VAE)를 이해하기 전에, 먼저 일반적인 오토인코더(Autoencoder, AE)의 작동 방식과 그 한계를 살펴보겠습니다.

\begin{tcolorbox}[conceptbox, title=일반적인 오토인코더 (AE)]
\textbf{한 줄 요약:} 입력 데이터를 압축(인코딩)했다가 다시 복원(디코딩)하는 신경망.

\textbf{작동 원리:}
\begin{enumerate}
    \item \textbf{인코더 (Encoder):} 입력 데이터 $X$ (예: 이미지)를 받아 저차원의 "잠재 벡터(Latent Vector)" $C$ (또는 인코딩)로 압축합니다.
    \item \textbf{디코더 (Decoder):} 잠재 벡터 $C$를 받아 원래 입력 데이터 $X$와 유사한 출력 $X'$를 재구성합니다.
\end{enumerate}
\textbf{목표:} 입력 $X$와 출력 $X'$ 사이의 재구성 손실(Reconstruction Loss)을 최소화하여, 잠재 벡터 $C$가 $X$의 효율적인 표현이 되도록 학습합니다.
\end{tcolorbox}

\begin{tcolorbox}[warningbox, title=오토인코더의 잠재 공간 문제: "구멍(Holes)"]
일반적인 오토인코더는 잠재 공간(Latent Space)에 문제가 있습니다.
\begin{itemize}
    \item \textbf{문제점:} 훈련된 오토인코더에서 특정 입력 $X$에 대응하는 잠재 벡터 $C$는 실제 데이터를 잘 표현합니다. 하지만 이 $C$를 잠재 공간에서 아주 조금만 이동시켜 새로운 $C'$를 만들고 디코더에 넣으면, 결과 $X''$는 전혀 현실적이지 않거나 의미 없는 데이터가 나올 수 있습니다.
    \item \textbf{원인:} 일반적인 오토인코더는 잠재 공간을 "구조화"하도록 강제되지 않습니다. 즉, 훈련 데이터에 대응하는 잠재 벡터들은 잠재 공간에 띄엄띄엄 존재하며, 그 사이의 공간(구멍)은 어떤 의미 있는 데이터에도 대응하지 않을 수 있습니다.
    \item \textbf{결과:} 잠재 공간에서 연속적인 변화를 통해 현실적인 데이터를 생성하거나, 한 이미지에서 다른 이미지로 부드럽게 변환하는 등의 작업이 불가능합니다.
\end{itemize}
\end{tcolorbox}

\subsubsection{변분 오토인코더 (VAE)의 해결책}
VAE는 일반 오토인코더의 잠재 공간 문제를 해결하기 위해 "확률적 인코딩"과 "잠재 손실(Latent Loss)" 개념을 도입합니다.

\begin{tcolorbox}[conceptbox, title=변분 오토인코더 (VAE)의 핵심 아이디어]
\textbf{한 줄 요약:} 잠재 공간에 "무작위성"을 주입하여, 그 공간을 연속적이고 의미 있게 만든다.

\textbf{직관적 예시:} 일반 오토인코더가 특정 그림을 그릴 때, 정확히 정해진 색상 팔레트의 한 점만 사용한다면, VAE는 그 점 주변의 다양한 색상들을 무작위로 조금씩 섞어가며 그림을 그립니다. 이렇게 하면 팔레트의 모든 점들이 어떤 색상에든 대응하게 되어, 팔레트 위를 부드럽게 이동하며 색상을 바꿀 수 있게 됩니다.

\textbf{기술적 설명:}
\begin{enumerate}
    \item \textbf{확률적 인코더:} VAE의 인코더는 입력 $X$를 잠재 벡터 $C$로 직접 변환하는 대신, 잠재 공간의 \textbf{확률 분포의 파라미터}를 출력합니다. 일반적으로 각 잠재 차원에 대해 평균($\mu$)과 분산($\sigma^2$)을 출력하여 정규 분포($\mathcal{N}(\mu, \sigma^2)$)를 정의합니다.
    \item \textbf{잠재 벡터 샘플링:} 인코더가 정의한 분포에서 잠재 벡터 $C$를 \textbf{무작위로 샘플링}합니다. 이 무작위성은 VAE의 핵심입니다.
    \item \textbf{디코더:} 샘플링된 잠재 벡터 $C$를 받아 $X'$를 재구성합니다.
\end{enumerate}
\end{tcolorbox}

\subsubsection{VAE의 손실 함수}
VAE는 두 가지 종류의 손실 함수를 최소화하여 훈련됩니다.
\begin{enumerate}
    \item \textbf{재구성 손실 (Reconstruction Loss):} 일반 오토인코더와 동일하게, 입력 $X$와 재구성된 $X'$ 사이의 차이를 측정합니다. 이는 모델이 원본 데이터를 잘 복원하도록 합니다.
    \item \textbf{KL 발산 손실 (KL Divergence Loss):} 인코더가 출력하는 잠재 분포($\mathcal{N}(\mu, \sigma^2)$)가 미리 정의된 \textbf{표준 정규 분포($\mathcal{N}(0, 1)$)}와 얼마나 다른지를 측정합니다.
    \[ \mathcal{L}_{KL} = D_{KL}(\mathcal{N}(\mu, \sigma^2) || \mathcal{N}(0, 1)) \]
    \begin{itemize}
        \item \textbf{필요성:} 만약 KL 발산 손실이 없다면, 모델은 재구성 손실을 최소화하기 위해 $\sigma$를 0으로 수렴시키려 할 것입니다. 이는 잠재 벡터 $C$를 샘플링하는 대신 고정된 값으로 만들어서 일반 오토인코더와 동일한 문제가 발생합니다.
        \item \textbf{효과:} KL 발산 손실은 모든 잠재 분포의 평균($\mu$)을 0에 가깝게 만들고, 분산($\sigma^2$)을 1에 가깝게 만들도록 강제합니다. 이는 잠재 공간의 모든 지점이 의미 있는 데이터에 대응하도록 하여 "구조화된" 잠재 공간을 만듭니다.
    \end{itemize}
\end{enumerate}
\textbf{총 손실:} $\mathcal{L}_{total} = \mathcal{L}_{reconstruction} + \mathcal{L}_{KL}$

\subsubsection{VAE의 장점과 한계}
\begin{itemize}
    \item \textbf{장점:}
    \begin{itemize}
        \item \textbf{구조화된 잠재 공간:} 잠재 공간이 연속적이고 부드러워, 잠재 벡터를 조작하여 의미 있는 새로운 데이터를 생성하거나 데이터 간의 부드러운 변환(예: 웃는 얼굴에서 찡그린 얼굴로 변화)이 가능합니다.
        \item \textbf{생성 능력:} 잠재 공간에서 임의의 벡터를 샘플링하여 디코더에 넣으면 현실적인 새로운 데이터를 생성할 수 있습니다.
    \end{itemize}
    \item \textbf{한계:}
    \begin{itemize}
        \item \textbf{생성 품질:} 생성된 데이터의 품질이 GAN(Generative Adversarial Networks)만큼 뛰어나지 않을 수 있습니다. 무작위성 때문에 재구성된 출력이 원본보다 약간 "흐릿"하거나 덜 선명할 수 있습니다.
        \item \textbf{훈련 복잡성:} 두 가지 손실 항을 균형 있게 최적화해야 합니다.
    \end{itemize}
\end{itemize}

\newpage
\subsection{생성적 적대 신경망 (Generative Adversarial Networks, GAN)}

\subsubsection{개념 및 필요성}
생성적 적대 신경망(GAN)은 지난 20년간 신경망 분야에서 가장 중요한 발전 중 하나로 꼽히는 혁신적인 생성 모델입니다. VAE와 마찬가지로 새로운 데이터를 생성하는 것이 목표지만, 완전히 다른 접근 방식을 사용합니다.

\begin{tcolorbox}[conceptbox, title=생성적 적대 신경망 (GAN)의 핵심 아이디어]
\textbf{한 줄 요약:} 두 개의 신경망이 서로 경쟁하며, 하나는 가짜를 만들고 다른 하나는 진짜를 구별하는 법을 배운다.

\textbf{직관적 예시:}
\begin{itemize}
    \item \textbf{위조지폐범 (생성자):} 실제 지폐와 구별하기 어려운 가짜 지폐를 만들려고 끊임없이 노력합니다.
    \item \textbf{경찰 (판별자):} 위조지폐범이 만든 지폐가 진짜인지 가짜인지 구별하는 능력을 키웁니다.
\end{itemize}
이 둘은 서로를 속이고 구별하려는 경쟁을 통해, 위조지폐범은 점점 더 정교한 가짜 지폐를 만들게 되고, 경찰은 점점 더 날카롭게 위조지폐를 구별하게 됩니다. 결국, 위조지폐범은 경찰도 속일 수 있는 완벽한 가짜 지폐를 만들게 됩니다.
\end{itemize}

\textbf{기술적 설명:} GAN은 두 개의 신경망으로 구성됩니다.
\begin{enumerate}
    \item \textbf{생성자 (Generator, G):} 무작위 노이즈(random noise)를 입력으로 받아 실제 데이터와 유사한 가짜 데이터(예: 가짜 이미지)를 생성합니다. 생성자의 목표는 판별자를 속이는 것입니다.
    \item \textbf{판별자 (Discriminator, D):} 실제 데이터와 생성자가 만든 가짜 데이터를 입력으로 받아, 입력된 데이터가 실제인지 가짜인지 이진 분류(binary classification)합니다. 판별자의 목표는 생성자를 속이지 않고 정확하게 구별하는 것입니다.
\end{enumerate}
이 두 네트워크는 "제로섬 게임(zero-sum game)" 방식으로 서로 경쟁하며 동시에 훈련됩니다.
\end{tcolorbox}

\subsubsection{GAN의 훈련 과정}
GAN의 훈련은 다음과 같은 반복적인 과정으로 이루어집니다.
\begin{enumerate}
    \item \textbf{판별자 훈련:}
    \begin{itemize}
        \item 실제 데이터를 입력으로 받아 '진짜(real)'라고 예측하도록 훈련합니다.
        \item 생성자가 만든 가짜 데이터를 입력으로 받아 '가짜(fake)'라고 예측하도록 훈련합니다.
        \item 판별자는 실제와 가짜를 정확히 구별하는 능력을 최대화하는 방향으로 학습합니다.
    \end{itemize}
    \item \textbf{생성자 훈련:}
    \begin{itemize}
        \item 무작위 노이즈를 입력으로 받아 가짜 데이터를 생성합니다.
        \item 생성된 가짜 데이터를 판별자에 입력합니다.
        \item 생성자는 판별자가 자신의 데이터를 '진짜'라고 예측하도록 속이는 방향으로 학습합니다. 즉, 판별자의 손실을 최소화하는 것이 아니라, 판별자가 '가짜'를 '진짜'라고 판단할 확률을 최대화하는 방향으로 학습합니다.
    \end{itemize}
\end{enumerate}
이 과정이 충분히 반복되면, 생성자는 실제 데이터와 거의 구별할 수 없는 데이터를 생성하게 되고, 판별자는 더 이상 실제와 가짜를 구별하기 어려워집니다 (즉, 50\% 확률로 예측).

\subsubsection{GAN의 장점과 한계}
\begin{itemize}
    \item \textbf{장점:}
    \begin{itemize}
        \item \textbf{비지도 학습:} 레이블이 없는 대량의 데이터만으로도 매우 현실적이고 고품질의 데이터를 생성할 수 있습니다.
        \item \textbf{뛰어난 생성 품질:} 특히 이미지 생성 분야에서 VAE를 능가하는 놀라운 현실감을 보여줍니다. 존재하지 않는 사람의 얼굴이나 고양이 이미지를 생성하는 것이 가능합니다.
        \item \textbf{다양한 응용:} 이미지 생성, 스타일 변환, 해상도 향상, 데이터 증강 등 다양한 분야에 적용됩니다.
    \end{itemize}
    \item \textbf{한계:}
    \begin{itemize}
        \item \textbf{훈련의 불안정성:} GAN은 훈련하기 매우 어렵고 불안정합니다. 생성자와 판별자 간의 균형을 맞추기 어려워, 한쪽이 너무 강해지거나 "모드 붕괴(Mode Collapse)"(생성자가 제한된 종류의 데이터만 생성하는 현상)와 같은 문제가 발생할 수 있습니다.
        \item \textbf{게임 이론적 문제:} 두 네트워크가 서로 다른 목표를 가지고 경쟁하기 때문에, 최적의 균형점을 찾는 것이 복잡합니다.
    \end{itemize}
\end{itemize}

\begin{tcolorbox}[warningbox, title=GAN 훈련의 불안정성 해결 노력]
강의에서 언급된 바와 같이, GAN 훈련의 불안정성은 게임 이론의 "비협력 게임" 문제와 유사합니다. 한 네트워크가 특정 "속임수"를 학습하면 다른 네트워크가 이에 적응하고, 이로 인해 전체 시스템이 특정 유형의 데이터에만 집중하거나 학습이 정체될 수 있습니다.
이를 해결하기 위해 \textbf{경험 재생(Experience Replay)}과 같은 기법이 사용됩니다. 생성자가 과거에 생성했던 이미지들을 다시 판별자에게 보여줌으로써, 판별자가 특정 유형의 가짜 데이터에만 특화되어 다른 유형을 잊어버리는 것을 방지하고, 더 일반적인 구별 능력을 유지하도록 돕습니다. StyleGAN과 같은 고급 GAN 아키텍처는 이러한 훈련 안정성과 생성 품질을 크게 향상시켰습니다.
\end{tcolorbox}

\newpage
\section*{절차/방법}

각 모델의 훈련 및 사용 절차를 간략히 요약합니다.

\subsection{마르코프 체인 및 은닉 마르코프 모델 (HMM)}
\begin{enumerate}
    \item \textbf{모델 정의:} 상태 공간(관찰/은닉), 전이 확률, 방출 확률, 초기 상태 확률을 정의합니다.
    \item \textbf{파라미터 추정 (훈련):}
    \begin{itemize}
        \item \textbf{은닉 상태가 관찰 가능한 경우:} 훈련 데이터에서 각 상태 전이 및 방출의 빈도수를 세어 확률을 직접 계산합니다 (최대 우도 추정).
        \item \textbf{은닉 상태가 관찰 불가능한 경우:} 기대-최대화(EM) 알고리즘을 사용하여 파라미터를 반복적으로 추정합니다.
    \end{itemize}
    \item \textbf{사용 (예측):}
    \begin{itemize}
        \item \textbf{가장 가능성 높은 은닉 상태 시퀀스 찾기:} 새로운 관찰 시퀀스($X$)가 주어졌을 때, 훈련된 HMM을 사용하여 가장 가능성 높은 은닉 상태 시퀀스($Y$)를 찾습니다 (예: 비터비 알고리즘).
        \item \textbf{관찰 시퀀스의 확률 계산:} 특정 관찰 시퀀스가 모델에서 생성될 확률을 계산합니다 (예: 포워드 알고리즘).
    \end{itemize}
\end{enumerate}

\subsection{조건부 무작위 필드 (CRF)}
\begin{enumerate}
    \item \textbf{특징 함수 설계:} 도메인 지식을 바탕으로 입력 시퀀스 $X$와 출력 레이블 $Y$ 간의 관계를 포착하는 특징 함수 $f_k$를 수동으로 정의합니다.
    \item \textbf{파라미터 추정 (훈련):} 레이블이 지정된 훈련 데이터($X, Y$)를 사용하여 각 특징 함수에 대한 가중치 $\lambda_k$를 최대 우도 추정(Maximum Likelihood Estimation)을 통해 학습합니다. 일반적으로 경사 하강법(Gradient Descent) 기반의 최적화 알고리즘이 사용됩니다.
    \item \textbf{사용 (예측):} 새로운 입력 시퀀스 $X$가 주어졌을 때, 훈련된 CRF 모델을 사용하여 가장 가능성 높은 출력 레이블 시퀀스 $Y$를 예측합니다 (예: 비터비 알고리즘).
\end{enumerate}

\subsection{양방향 LSTM과 CRF 결합 (Bi-LSTM-CRF)}
\begin{enumerate}
    \item \textbf{모델 아키텍처 구축:} 임베딩 레이어, 양방향 LSTM 레이어, 선형 변환 레이어, CRF 레이어를 순서대로 연결하여 모델을 구성합니다.
    \item \textbf{훈련 데이터 준비:} 입력 문장($X$)을 단어 인덱스 시퀀스로, 출력 레이블($Y$)을 태그 인덱스 시퀀스로 변환합니다.
    \item \textbf{모델 훈련:}
    \begin{itemize}
        \item 입력 문장 시퀀스를 임베딩 레이어에 통과시켜 단어 임베딩을 얻습니다.
        \item 임베딩을 양방향 LSTM에 통과시켜 각 단어에 대한 문맥적 은닉 상태를 얻습니다.
        \item 은닉 상태를 선형 변환 레이어에 통과시켜 각 태그에 대한 방출 점수(emission scores)를 계산합니다.
        \item 방출 점수와 실제 태그 시퀀스를 CRF 레이어에 전달하여 손실(음의 로그 우도)을 계산합니다.
        \item 경사 하강법을 사용하여 모델의 모든 파라미터(임베딩, LSTM, 선형 변환, CRF 전이 점수)를 종단 간(end-to-end)으로 최적화합니다.
    \end{itemize}
    \item \textbf{사용 (예측):} 새로운 입력 문장이 주어졌을 때, 훈련된 모델의 \texttt{decode} 메서드를 사용하여 가장 가능성 높은 태그 시퀀스를 예측합니다.
\end{enumerate}

\subsection{변분 오토인코더 (VAE)}
\begin{enumerate}
    \item \textbf{모델 아키텍처 구축:} 인코더(입력 $X$를 $\mu$와 $\log(\sigma^2)$로 매핑), 샘플링 레이어(재매개변수화 트릭 사용), 디코더(샘플링된 잠재 벡터 $C$를 $X'$로 매핑)를 구성합니다.
    \item \textbf{훈련 데이터 준비:} 레이블이 없는 원본 데이터 $X$를 준비합니다.
    \item \textbf{모델 훈련:}
    \begin{itemize}
        \item 입력 $X$를 인코더에 통과시켜 잠재 분포의 평균 $\mu$와 로그 분산 $\log(\sigma^2)$를 얻습니다.
        \item $\mu$와 $\log(\sigma^2)$를 사용하여 잠재 벡터 $C$를 샘플링합니다 (재매개변수화 트릭을 통해 역전파 가능하게 함).
        \item 샘플링된 $C$를 디코더에 통과시켜 재구성된 $X'$를 얻습니다.
        \item \textbf{손실 함수 계산:} 재구성 손실(예: MSE 또는 BCE)과 KL 발산 손실을 합산하여 총 손실을 계산합니다.
        \item 경사 하강법을 사용하여 인코더와 디코더의 모든 파라미터를 최적화합니다.
    \end{itemize}
    \item \textbf{사용 (생성):}
    \begin{itemize}
        \item \textbf{새로운 데이터 생성:} 표준 정규 분포($\mathcal{N}(0, 1)$)에서 임의의 잠재 벡터 $C$를 샘플링하여 디코더에 입력하면 새로운 데이터를 생성할 수 있습니다.
        \item \textbf{데이터 변환/보간:} 두 데이터 포인트의 잠재 벡터 사이를 선형 보간하여 부드러운 변환 효과를 얻을 수 있습니다.
    \end{itemize}
\end{enumerate}

\subsection{생성적 적대 신경망 (GAN)}
\begin{enumerate}
    \item \textbf{모델 아키텍처 구축:} 생성자(Generator, G)와 판별자(Discriminator, D) 두 개의 신경망을 독립적으로 구축합니다.
    \item \textbf{훈련 데이터 준비:} 레이블이 없는 실제 데이터 $X_{real}$를 준비합니다.
    \item \textbf{모델 훈련 (반복적):}
    \begin{itemize}
        \item \textbf{판별자 훈련 단계:}
        \begin{itemize}
            \item 실제 데이터 $X_{real}$를 D에 입력하고, D가 '진짜'라고 예측하도록 손실을 계산하고 D의 파라미터를 업데이트합니다.
            \item 무작위 노이즈 $Z$를 G에 입력하여 가짜 데이터 $X_{fake}$를 생성합니다.
            \item $X_{fake}$를 D에 입력하고, D가 '가짜'라고 예측하도록 손실을 계산하고 D의 파라미터를 업데이트합니다.
        \end{itemize}
        \item \textbf{생성자 훈련 단계:}
        \begin{itemize}
            \item 무작위 노이즈 $Z$를 G에 입력하여 가짜 데이터 $X_{fake}$를 생성합니다.
            \item $X_{fake}$를 D에 입력하고, D가 이 데이터를 '진짜'라고 예측하도록 손실을 계산하고 G의 파라미터를 업데이트합니다 (이때 D의 파라미터는 고정).
        \end{itemize}
        \item 이 두 단계를 번갈아 가며 반복하여 G와 D가 서로 경쟁하며 학습하도록 합니다.
    \end{itemize}
    \item \textbf{사용 (생성):} 훈련된 생성자 G에 무작위 노이즈를 입력하여 새로운, 현실적인 데이터를 생성합니다.
\end{enumerate}

\newpage
\section*{체크리스트}

이 체크리스트는 강의 내용을 효과적으로 학습하고 이해했는지 점검하는 데 도움이 됩니다.

\begin{tcolorbox}[summarybox, title=학습 점검 체크리스트]
\begin{itemize}
    \item \textbf{퀴즈 10 복습}
    \begin{itemize}
        \item NER의 목적, 규칙 기반 NER, spaCy의 역할, 통계적 방법의 단점을 설명할 수 있는가?
        \item spaCy의 시각화 기능이 GUI가 아닌 이유를 이해하는가?
    \end{itemize}
    \item \textbf{마르코프 체인 (Markov Chains)}
    \begin{itemize}
        \item 마르코프 속성을 정의하고 예시를 들 수 있는가?
        \item 마르코프 속성의 한계와 "상태" 재정의를 통한 극복 방법을 설명할 수 있는가?
    \end{itemize}
    \item \textbf{은닉 마르코프 모델 (HMM)}
    \begin{itemize}
        \item HMM이 필요한 이유와 은닉 상태/관찰 상태의 개념을 이해하는가?
        \item 전이 확률, 방출 확률, 초기 상태 확률을 정의할 수 있는가?
        \item HMM의 훈련 방법(EM 알고리즘)과 주요 적용 사례를 아는가?
        \item HMM의 주요 한계(단순화된 의존성)를 설명할 수 있는가?
    \end{itemize}
    \item \textbf{조건부 무작위 필드 (CRF)}
    \begin{itemize}
        \item CRF가 HMM의 어떤 한계를 극복하는지 설명할 수 있는가?
        \item 특징 함수($f_k$)의 역할과 수동 특징 공학의 어려움을 이해하는가?
        \item CRF의 장점(유연성, 문맥 인식)과 단점(집중적인 특징 공학)을 설명할 수 있는가?
    \end{itemize}
    \item \textbf{양방향 LSTM과 CRF 결합 (Bi-LSTM-CRF)}
    \begin{itemize}
        \item Bi-LSTM-CRF가 CRF의 특징 공학 문제를 어떻게 해결하는지 설명할 수 있는가?
        \item 모델의 아키텍처(임베딩, Bi-LSTM, 선형 변환, CRF)와 각 구성 요소의 역할을 이해하는가?
        \item "방출 점수(Emission Scores)"의 의미와 역할은 무엇인가?
        \item Bi-LSTM-CRF의 장점(자동 특징 학습)과 단점(높은 계산 비용)을 아는가?
    \end{itemize}
    \item \textbf{변분 오토인코더 (VAE)}
    \begin{itemize}
        \item 일반 오토인코더의 잠재 공간 문제("구멍")를 설명할 수 있는가?
        \item VAE가 잠재 공간을 "구조화"하는 방법을 이해하는가(확률적 인코더, 샘플링)?
        \item VAE의 손실 함수(재구성 손실 + KL 발산 손실)와 KL 발산의 중요성을 설명할 수 있는가?
        \item VAE의 장점(구조화된 잠재 공간, 생성 능력)과 한계(생성 품질)를 아는가?
    \end{itemize}
    \item \textbf{생성적 적대 신경망 (GAN)}
    \begin{itemize}
        \item GAN이 필요한 이유와 생성자/판별자의 역할을 설명할 수 있는가?
        \item GAN의 훈련 과정(경쟁적 학습)을 이해하는가?
        \item GAN의 장점(뛰어난 생성 품질, 비지도 학습)과 한계(훈련 불안정성, 모드 붕괴)를 아는가?
        \item GAN 훈련의 불안정성 해결 노력(예: 경험 재생)을 이해하는가?
    \end{itemize}
\end{itemize}
\end{tcolorbox}

\newpage
\section*{자주 묻는 질문 (FAQ)}

\begin{tcolorbox}[conceptbox, title=Q\&A: 초심자를 위한 핵심 질문]
\textbf{Q1: 마르코프 체인, HMM, CRF는 어떻게 다른가요?}
\textbf{A1:} 마르코프 체인은 모든 상태가 관찰 가능하고 미래가 현재에만 의존하는 가장 기본적인 모델입니다. HMM은 관찰 불가능한 은닉 상태를 도입하여 마르코프 체인을 확장했지만, 여전히 강력한 독립성 가정을 가집니다. CRF는 HMM의 독립성 가정을 완화하여 입력 시퀀스 전체의 문맥을 고려할 수 있는 훨씬 유연한 모델입니다.

\textbf{Q2: CRF의 "특징 함수"를 수동으로 만드는 것이 왜 그렇게 어려운가요?}
\textbf{A2:} 특징 함수는 도메인 전문가가 특정 패턴(예: "Mr." 뒤에 대문자 단어)을 직접 규칙으로 코딩하는 것입니다. 현실 세계의 언어나 이미지 패턴은 매우 복잡하고 다양하기 때문에, 모든 중요한 패턴을 수동으로 찾아내고 규칙으로 만드는 것은 엄청난 시간과 노력이 필요하며, 놓치는 부분이 많을 수 있습니다.

\textbf{Q3: Bi-LSTM-CRF에서 Bi-LSTM과 CRF는 각각 어떤 역할을 하나요?}
\textbf{A3:} Bi-LSTM은 입력 시퀀스(예: 문장)의 과거와 미래 문맥 정보를 모두 활용하여 각 단어에 대한 풍부한 "특징"을 자동으로 추출합니다. CRF는 Bi-LSTM이 추출한 이 특징들을 바탕으로, 출력 레이블 시퀀스(예: 품사 태그) 간의 논리적이고 유효한 전이 규칙을 적용하여 최종적으로 가장 적합한 레이블 시퀀스를 결정합니다.

\textbf{Q4: 일반 오토인코더와 변분 오토인코더(VAE)의 가장 큰 차이점은 무엇인가요?}
\textbf{A4:} 일반 오토인코더는 입력 데이터를 하나의 고정된 잠재 벡터로 압축합니다. 이로 인해 잠재 공간에 "구멍"이 생겨 의미 있는 데이터 생성이 어렵습니다. VAE는 입력 데이터를 잠재 공간의 "확률 분포"로 인코딩하고, 이 분포에서 잠재 벡터를 샘플링합니다. KL 발산 손실을 통해 잠재 공간을 구조화하여 연속적이고 의미 있는 데이터 생성을 가능하게 합니다.

\textbf{Q5: GAN은 왜 훈련하기가 그렇게 어려운가요?}
\textbf{A5:} GAN은 생성자와 판별자라는 두 네트워크가 서로 경쟁하는 "게임" 형태로 훈련됩니다. 이 게임의 균형을 맞추기가 매우 어렵습니다. 한쪽이 너무 빨리 강해지면 다른 쪽이 제대로 학습하지 못하거나, 생성자가 특정 종류의 데이터만 반복해서 만드는 "모드 붕괴" 현상이 발생할 수 있습니다. 이는 안정적인 학습을 위한 섬세한 튜닝을 요구합니다.
\end{tcolorbox}

\newpage
\section*{빠르게 훑어보기 (1페이지 요약)}

\begin{tcolorbox}[summarybox, title=핵심 개념 요약]
\textbf{1. 마르코프 체인 (Markov Chain)}
\begin{itemize}
    \item \textbf{핵심:} 미래는 현재에만 의존 (마르코프 속성).
    \item \textbf{장점:} 간단한 시스템 동역학 모델링.
    \item \textbf{단점:} 강력한 독립성 가정, 복잡한 문맥 무시.
    \item \textbf{해결:} "상태" 정의를 확장 (예: 2-gram).
\end{itemize}
\end{tcolorbox}

\begin{tcolorbox}[summarybox, title=핵심 개념 요약]
\textbf{2. 은닉 마르코프 모델 (HMM)}
\begin{itemize}
    \item \textbf{핵심:} 관찰 불가능한 은닉 상태($Y$)가 관찰 가능한 상태($X$)를 생성.
    \item \textbf{파라미터:} 전이 확률($Y \to Y$), 방출 확률($Y \to X$), 초기 상태 확률.
    \item \textbf{훈련:} EM 알고리즘 (은닉 상태가 관찰 불가능할 때).
    \item \textbf{장점:} 시퀀스 데이터 모델링, 품사 태깅, 음성 인식.
    \item \textbf{단점:} 여전히 강력한 독립성 가정, 문맥 정보 활용 제한.
\end{itemize}
\end{tcolorbox}

\begin{tcolorbox}[summarybox, title=핵심 개념 요약]
\textbf{3. 조건부 무작위 필드 (CRF)}
\begin{itemize}
    \item \textbf{핵심:} HMM의 독립성 한계 극복, 입력 시퀀스 전체($X$)를 고려하여 출력 시퀀스($Y$) 예측.
    \item \textbf{구성:} 특징 함수($f_k$)와 가중치($\lambda_k$).
    \item \textbf{장점:} 유연한 문맥 포착, 장거리 의존성 모델링.
    \item \textbf{단점:} \textbf{수동 특징 공학} 필요 (가장 큰 문제점).
\end{itemize}
\end{tcolorbox}

\begin{tcolorbox}[summarybox, title=핵심 개념 요약]
\textbf{4. 양방향 LSTM과 CRF 결합 (Bi-LSTM-CRF)}
\begin{itemize}
    \item \textbf{핵심:} Bi-LSTM이 자동으로 특징을 학습하고, CRF가 시퀀스 제약 조건을 적용.
    \item \textbf{작동:} Bi-LSTM $\to$ 은닉 상태 $\to$ 선형 변환 $\to$ 방출 점수 $\to$ CRF 레이어.
    \item \textbf{장점:} \textbf{자동 특징 학습}, 문맥 포착, 시퀀스 레이블링 성능 우수.
    \item \textbf{단점:} 높은 계산 비용, 복잡한 모델 튜닝.
\end{itemize}
\end{tcolorbox}

\begin{tcolorbox}[summarybox, title=핵심 개념 요약]
\textbf{5. 변분 오토인코더 (VAE)}
\begin{itemize}
    \item \textbf{핵심:} 잠재 공간을 확률 분포로 인코딩하여 "구조화"된 잠재 공간 생성.
    \item \textbf{구성:} 인코더($\mu, \sigma^2$ 출력), 샘플링, 디코더.
    \item \textbf{손실:} 재구성 손실 + KL 발산 손실 (잠재 분포를 표준 정규 분포에 가깝게).
    \item \textbf{장점:} 의미 있는 데이터 생성, 잠재 공간 보간 가능.
    \item \textbf{단점:} 생성 품질이 GAN보다 낮을 수 있음.
\end{itemize}
\end{tcolorbox}

\begin{tcolorbox}[summarybox, title=핵심 개념 요약]
\textbf{6. 생성적 적대 신경망 (GAN)}
\begin{itemize}
    \item \textbf{핵심:} 생성자(가짜 생성)와 판별자(진짜/가짜 구별)가 경쟁하며 학습.
    \item \textbf{훈련:} 제로섬 게임, 서로를 속이고 구별하며 발전.
    \item \textbf{장점:} \textbf{매우 현실적인 데이터 생성}, 비지도 학습.
    \item \textbf{단점:} 훈련 불안정성, 모드 붕괴, 균형 맞추기 어려움.
\end{itemize}
\end{tcolorbox}

\end{document}